\documentclass[10pt]{article}

\usepackage[utf8]{inputenc}

\usepackage[T1]{fontenc}

\usepackage{amsmath}

\usepackage{amsfonts}

\usepackage{amssymb}

\usepackage[version=4]{mhchem}

\usepackage{stmaryrd}

\usepackage{bbold}

\usepackage{graphicx}

\usepackage[export]{adjustbox}

\graphicspath{ {./images/} }



\begin{document}

щее, что между значениями Ф. из данной системы существуют определенные связи, в частности значения одной из них выражаются через значения остальных. Hamp., Ф. $f_{1}(x)=\sin ^{2} x$ и $f_{2}(x)=\cos ^{2}(x)$ зависимы на всей числовой прямой, поскольку всегда $\sin ^{2} x=$ $=1-\cos ^{2} x$. Пусть $D$-область в $\mathbb{R} n, \bar{D}$-ее замыкание, $f: \bar{D} \longrightarrow \mathbb{R} n, f(x)=\left\{y_{i}=f_{i}(x), i=1,2, \ldots, n\right\}, \quad x \in \bar{D}$. Функции $f_{i}, i=1,2, \ldots, n, \quad$ наз. $з$ а в и с и м ы м и на $\bar{D}$, если существует такая неп рерывно дифференцируемая на $\mathbb{R}^{n}$ функция $F(y), y \in \mathbb{R} n$, множество нулей к-рой образует нигде не плотное в пространстве $\mathbb{R} n$ множество, qто композиция $F \circ f$ тождественно равна нулю на $\bar{D}$.



Функции $f_{i}: G \longrightarrow \mathbb{R}, i=1,2, \ldots, n$, наз. $з$ а в ис им м и на области $G \subset \mathbb{R} n$, если они зависимы на замыкании $\bar{D}$ любой области $D \subset G$ такой, что $\bar{D} \subset G$.



Непрерывно дифференцируемые в области $G \subset \mathbb{R}^{n}$ функции $f_{i}, i=1,2, \ldots, n$, зависимы на $G$ тогда и только тогда, когда их якобиан



\[

\frac{\partial\left(f_{1}, \ldots, f_{n}\right)}{\partial\left(x_{1}, \ldots, x_{n}\right)}

\]



тождественно равен нулю на области $G$.



Пусть теперь



\[

f_{i}: G \longrightarrow \mathbb{R}^{m}, \quad i=1,2, \ldots, m \leqslant n, G \subset \mathbb{R}^{n} .

\]



Если для Ф. $y_{i}=f_{i}(x), \quad x \in G, \quad i=1,2, \ldots, m$, $G \subset \mathbb{R}^{n}$, существуют открытое в $R_{y_{1} y_{2}}^{m-1} \ldots y_{m-1}$ множество $\Gamma$ и непрерывно диффференцируемая на $\Gamma$ функция $\Phi\left(y_{1}, \ldots, y_{m-1}\right)$ такие, что в любой точке $x \in G$ выполняются условия



\[

\left(f_{1}(x), \ldots, f_{m-1}(x)\right) \in \Gamma

\]



n



\[

\Phi\left(f_{1}(x), \ldots, f_{m-1}(x)\right)=f_{m}(x) \text {, }

\]



то Ф. $f_{m}$ наз. за ви с им ой на множе с тве $G$ от Ф. $f_{1}, \ldots, f_{m-1}$.



Если Ф. $f_{i}, i=1,2, \ldots, m=n$, кепрерывны в области $G$ И в окрестности каждой точки $x \in G$ одна из них зависит от остальных, то Ф. $f_{i}, i=1,2, \ldots, n$, зависимы на области $G$.



Если в окрестности каждоӥ точки $x \in G$ одна из непрерывно дифференцируемых в области $G \subset \mathbb{R} n$ функций $f_{i}, i=1,2, \ldots, m \leqslant n$, зависит от остальных, то в любой точке области $G$ ранг матрицы Якоби



\[

\left\|\frac{\partial f_{i}}{\partial x_{j}}\right\|, \quad i=1,2, \ldots, m, \quad j=1,2, \ldots, n,

\]



меньше $m$, т. е. градиенты $\nabla f_{1}, \nabla f_{2}, \ldots, \nabla f_{m}$ линейно зависимы в каждой точке $x \in G$.



Пусть Ф. (1) непрерывно дифференцируемы на об ласти $G \subset \mathbb{R}^{n}$ и ранг их матрицы Якоби (2) в каждой точке $x \in G$ не превышает нек-рого числа $r, 1 \leqslant r<$ $<m \leqslant n$, а в нек-рой точке $x^{(0)} \in G$ равен $r$. Иначе говоря, существуют такие переменные $x_{j_{1}}, \ldots, x_{j_{r}}$ и функции $y_{i_{1}}=f_{i_{1}}(x), y_{i_{2}}=f_{i_{2}}(x), \ldots, y_{i_{r}}=f_{i_{r}}(x)$, что



\[

\left.\frac{\partial\left(f_{i_{1}}, f_{i_{2}}, \ldots, f_{i_{r}}\right)}{\partial\left(x_{j_{1}}, x_{j_{2}}, \cdot x_{j_{r}}\right)}\right|_{x=x^{(0)}} \neq 0 .

\]



Тогда ни в какой окрестности точки $x^{(0)}$ ни одна из Ф. $f_{i_{1}}, f_{i_{2}}, \ldots, f_{i_{r}}$ не является зависимой от остальных из них и существует такая окрестность точки $x^{(0)}$, что каждая из оставшихся $\Phi . f_{i}, i \neq i_{k}, k=1,2, \ldots, r$, зависит на әтой окрестности от $\Phi$. $f_{i_{1}}, f_{i_{2}}, \ldots, f_{i_{r}}$. В частности, если градиенты $\nabla f_{1}, \nabla f_{2}, \ldots, \nabla f_{m}$ линейно зависимы во всех точках области $G$ и в нек-рой точке $x^{(0)} \in G$ среди них $(m-1)$ линейно независимы, и, следовательно, один из них, напр. $\nabla f_{m}$, является линейной комбинацией остальных, то существует такая окрестность точки $x^{(0)}$, что на этой окрестности Ф. $f_{m}$ зависит от остальных Ф. $f_{1}, f_{2}, \ldots, f_{m-1}$.



Лит.: [1] Ә й л е p Л., Введение в анализ бесконечных, пер с лат., 2 изд., т. 1, М, 1961, [2] Л о б а ч е в с к и й Н. И., Полное собр. соч., т. 5, М.- Л., 1951; [3] Д е Д е к и н Д Р ,

Что такое числа и для чего они служат?, пер. с нем., Казань, 1905 , [4] X a у с д о р Ф Ф., Теория множеств, пер. с нем., М.-ЈI., 1937, [5] Т а р с к и й А., Введение в логику и методологию дедуктивных наук, пер. с англ., М., $1948,[6]$ к у р а т о вс к и й K., Топология, пер. с ангJ , Т. 1, M., 1966, 77] F r e g e G

Function, Begriff, Bedeutung, Göttingen, 1962, S. 79-95; [8] ч е р ч А., Введение в математическую люгику, пер. с англ., М., 1960, ' [9] ІО ш к е в и ч А. ПІ «Истор.-матем. исследования», 1966, в. 17, с 123-50; [10] М е д в е д е в Ф. А., Очерки истории теории функций действительного переменного, М., 1975.



\includegraphics[max width=\textwidth, center]{2023_07_09_fbe64bb64f9d0ddbc9dbg-1}

введения характерсв. L-ф. составляют сложной природы класс специальных функций комплексного переменного, определяемых Дирихле рядами или Эйлера произведениями с характерами. Они служат основным инструментом для изучения аналитич. методами арифметики соответствующих математич. объектов: поля рациональных чисел, алгебраич. полей, алгебраич. многообразий над конечными полями и т. П. Простейшими представителями $L$-ф. являются Дирихле $L$ функции. Остальные $L$-ф. представляют собой более или менее близкие аналоги и обобщения этих $L-\phi$.



ФУ. Ф. Лаврuк. ФУНКЦИЯ МНОЖЕСТВ - отображение $f$ нек-рой совокупности $\Sigma$ подмножеств данного множества $X$ в другое множество, обычно в множество $\mathbb{R}$ действительных или $\mathbb{C}$ комплексных чисел. Важным классом Ф. м. являются а д д и т и в ны е ф у н к ц и м н $о$ ж е с т в, для к-рых



\[

\begin{gathered}

f\left(\mathrm{U}_{i=1}^{n} M_{i}\right)=\sum_{i=1}^{n} f\left(M_{i}\right), \\

M_{i} \in \Sigma, M_{i} \cap M_{j}=\varnothing, i \neq j,

\end{gathered}

\]



и $\sigma-\mathbf{a} д$ д т ивны е ф ункции множөс т в, к-рые удовлетворяют равенству (:) и для счетной совокупности множеств. Если $f$ принимает лишь неотрицательные значения, $R(\varnothing)=0$ и $\Sigma$ является $\sigma$-алгеброй, то такая функция наз. м е р о й.



Лum.: [1] Ка н то р о в и ч Л. В., А к и Јо в Г. 1I., Функциональный анализ, 2 изд., М., 1977. В. И. Соболев.

ФУРБЕ ИНТЕГРАЛ - континуальный аналог Фурье ряда. Для функции, заданной на конечном промежутке действительной оси, важное значение имеет представлөние ее рядом Фурье. Для функции $f(x)$, заданной на всей оси, аналогичную роль играет разложение $f$ в интеграл Фурье:



\[

f(x)=\int_{0}^{\infty}[A(\lambda) \cos \lambda x+B(\lambda) \sin \lambda x] d x

\]



где



\[

\begin{aligned}

& A=\frac{1}{\pi} \int_{-\infty}^{+\infty} f(\xi) \cos \lambda \xi d \xi d \xi, \\

& B=\frac{1}{\pi} \int_{-\infty}^{+\infty} f(\xi) \sin \lambda \xi d \xi .

\end{aligned}

\]



Разложение (1) можно формально строить в предположениях, обеспечивающих существование написанных интегралов. Оно справедливо, напр., для гладкой финитной функции $f(x)$. Имеется большое число признаков, обеспечивающих равенство (1) в том или ином смысле. Подстановка (2) в (1) дает т. н. и н т еГральн ую форм у л у Ф у р ь е



\[

f(x)=\frac{1}{\pi} \int_{0}^{\infty} d \lambda \int_{-\infty}^{+\infty} f(\xi) \cos \lambda(x-\xi) d \xi,

\]



обоснование к-рой и приводит к упомянутым признакам. Большую пользу приносит при этом представление $f(x)$ п р о с тым ин те г ра лом Ф у р ье



\[

f(x)=\lim _{N \rightarrow \infty} \frac{1}{\pi} \int_{-\infty}^{+\infty} f(\xi) \frac{\sin N(x-\xi)}{x-\xi} d \xi,

\]





\end{document}